\documentclass[conference]{IEEEtran}
\usepackage{amsmath}
\usepackage{booktabs}
\usepackage{array}
\usepackage{url}

\title{Chest X-Ray Pneumonia Classification Using Deep Learning and Hybrid Feature Fusion}
\author{}

\begin{document}
\maketitle

\begin{abstract}
This report presents a deep learning and hybrid evaluation for binary classification of pediatric chest radiographs as normal or pneumonia. A five-block convolutional neural network (CNN) is trained from scratch, and a hybrid classifier combines its learned image representation with handcrafted intensity and edge features. The experiment uses a group-audited split of 3,168 training, 1,585 validation, and 618 test images. Three CPU training epochs are performed, with model selection based on validation balanced accuracy. On the test set, the CNN obtains 0.580 balanced accuracy and 0.995 pneumonia recall. The hybrid obtains 0.713 balanced accuracy and 0.984 recall. Despite the higher hybrid balanced accuracy, both systems produce substantial numbers of false positives. The evaluation is exploratory and does not establish clinical suitability.
\end{abstract}

\begin{IEEEkeywords}
chest radiography, pneumonia, convolutional neural networks, feature fusion, hybrid learning, image classification
\end{IEEEkeywords}

\section{Introduction}
Chest radiographs are commonly examined for signs associated with pneumonia. Image classification models can provide a useful educational setting for studying deep learning, but their performance depends on the data, split, and evaluation procedure. The objective of this project is to classify a radiograph as normal or pneumonia and examine whether a combination of learned and handcrafted features changes predictive performance.

The system includes a custom CNN trained from scratch, logistic regression applied to handcrafted features, logistic regression applied to CNN features, and a feature-fusion hybrid. The comparison evaluates both the deep learning model and the contribution of combining feature types. The work is a research exercise and is not a diagnostic tool.

\section{Dataset and Experimental Split}
The dataset contains pediatric chest X-ray images in normal and pneumonia categories. The image audit excludes exact pixel duplicates and groups identified as overlapping the test partition. The experiment uses 3,168 images for training, 1,585 for validation, and 618 for testing. Table~\ref{tab:split} summarizes class counts.

\begin{table}[ht]
\caption{Image counts in the experimental split}
\label{tab:split}
\centering
\begin{tabular}{lrrr}
\toprule
Split & Normal & Pneumonia & Total \\
\midrule
Training & 898 & 2,270 & 3,168 \\
Validation & 450 & 1,135 & 1,585 \\
Test & 231 & 387 & 618 \\
\bottomrule
\end{tabular}
\end{table}

Groups are based on filename and image-hash information rather than independently verified patient identities. The validation set is used for checkpoint and classifier hyperparameter selection. The test set is used for the reported comparison, but it has previously been inspected within the project; consequently, the results should be interpreted as exploratory.

\section{Data Preparation}
Images are read from disk in batches, avoiding the need to load the full dataset into memory. Each image is converted to grayscale, resized to $64\times64$ pixels, normalized to the range $[0,1]$, and repeated across three channels for the CNN. This resolution differs from the original notebook's $256\times256$ input and was chosen for the local CPU and memory budget. During training, random horizontal flipping provides augmentation. No augmentation is applied for validation or testing.

The labels are encoded as 0 for normal and 1 for pneumonia. Class-weighted cross-entropy gives greater weight to the smaller class during CNN training. This helps address class imbalance but does not guarantee high specificity or calibrated probabilities.

\section{CNN Architecture}
The CNN contains five feature-extraction blocks. Each block has two $3\times3$ convolutional layers with padding, batch normalization, rectified linear unit (ReLU) activations, max pooling, and dropout. The channel counts are 32, 64, 128, 256, and 512. With the $64\times64$ input used in this experiment, five pooling operations reduce the feature maps to $2\times2$. Adaptive average pooling then produces 512 values, followed by fully connected layers of sizes $512\rightarrow256\rightarrow128\rightarrow2$. The final outputs are two class logits.

For input $x$, the network computes an embedding $z=f_{\theta}(x)$ and class logits $g_{\theta}(z)$. Softmax transforms the logits into class probabilities. Batch normalization helps stabilize learning, while dropout with rate 0.3 and optimizer weight decay provide regularization.

\section{Training Procedure}
Training uses AdamW with learning rate 0.001, weight decay 0.001, class-weighted cross-entropy, and batch size 16. A random seed of 42 is used. The network is trained for three epochs on the local CPU. After each epoch, balanced accuracy is computed on the validation set; the checkpoint with the best validation balanced accuracy is retained for evaluation.

\begin{table}[ht]
\caption{CNN training history}
\label{tab:history}
\centering
\begin{tabular}{rrr}
\toprule
Epoch & Training loss & Validation balanced accuracy \\
\midrule
1 & 0.3005 & 0.5035 \\
2 & 0.2784 & 0.8278 \\
3 & 0.2284 & 0.5819 \\
\bottomrule
\end{tabular}
\end{table}

The epoch-two checkpoint is selected because it achieves the highest validation balanced accuracy. Training loss decreases across epochs, while validation balanced accuracy varies substantially. This variation suggests that more controlled training and repeated runs are needed before drawing conclusions about convergence.

\section{Hybrid Approach}
The hybrid uses two types of input for a logistic regression classifier. First, the trained CNN supplies a 512-dimensional representation after adaptive average pooling. Second, a handcrafted descriptor supplies 68 values: an $8\times8$ grid of average intensities, global mean and standard deviation, and average vertical and horizontal intensity differences. These values capture coarse image appearance and edge strength.

The two vectors are concatenated, standardized using training data, and passed to class-weighted logistic regression. The regularization parameter is selected from $C\in\{0.01,0.1,1.0\}$ according to validation balanced accuracy. Standardization is fitted on training data only. The hybrid is compared with logistic regression using only handcrafted features and only CNN features. This provides a feature-input ablation under the same data split and classifier family.

\section{Evaluation Metrics}
The evaluation uses accuracy, balanced accuracy, precision, recall, specificity, F1 score, and the area under the receiver operating characteristic curve (ROC-AUC). For true positives $TP$, true negatives $TN$, false positives $FP$, and false negatives $FN$,
\begin{align}
\mathrm{Recall} &= \frac{TP}{TP+FN}, \\
\mathrm{Specificity} &= \frac{TN}{TN+FP}, \\
\mathrm{Balanced\ Accuracy} &= \frac{\mathrm{Recall}+\mathrm{Specificity}}{2}.
\end{align}
Recall measures the proportion of pneumonia images identified, whereas specificity measures the proportion of normal images identified. ROC-AUC assesses ranking across thresholds; the other reported classification metrics use a threshold of 0.5.

\section{Results and Ablation}
Table~\ref{tab:results} reports test results for the four evaluated systems. The CNN head achieves high recall but low specificity, predicting many normal images as pneumonia. The hybrid obtains higher balanced accuracy than the CNN-feature-only logistic regression in this run. It also achieves higher balanced accuracy than the CNN head, although the handcrafted-only classifier performs similarly to the hybrid.

\begin{table}[ht]
\caption{Test-set results at threshold 0.5}
\label{tab:results}
\centering
\begin{tabular}{lrrrr}
\toprule
Model & Bal. acc. & Recall & Spec. & AUC \\
\midrule
CNN head & 0.580 & 0.995 & 0.165 & 0.847 \\
Handcrafted + LR & 0.715 & 0.979 & 0.450 & 0.875 \\
CNN features + LR & 0.685 & 0.941 & 0.429 & 0.866 \\
Hybrid fusion + LR & 0.713 & 0.984 & 0.442 & 0.882 \\
\bottomrule
\end{tabular}
\end{table}

The CNN head produces 385 true positives, 2 false negatives, 38 true negatives, and 193 false positives. The hybrid produces 381 true positives, 6 false negatives, 102 true negatives, and 129 false positives. Compared with the CNN head, the hybrid reduces false positives but misses four additional pneumonia images. The handcrafted-only model obtains 0.715 balanced accuracy, slightly above the hybrid's 0.713. Thus, the results do not demonstrate an advantage of fusion over the handcrafted features alone.

The feature ablation changes the classifier input while holding the split and logistic regression procedure fixed. It does not independently measure the impact of dropout, augmentation, image size, or individual CNN layers. No confidence intervals or repeated-seed experiments were computed for this five-block run.

\section{Limitations and Future Work}
The study uses one random seed, three epochs, and a compact image resolution. The dataset is limited to the available pediatric images, and the grouping procedure relies on filename and hash proxies rather than verified patient identifiers. Validation performance is higher than test performance; possible differences between the partitions and model overfitting require further investigation. Multiple models are reported on a test set that has previously been examined, so the results are exploratory.

Future work should repeat training with multiple seeds, verify patient-level grouping, evaluate independent external data, examine calibration and decision thresholds on validation data, and conduct targeted ablations of preprocessing and regularization. These steps are required before making broader claims about generalization or usefulness in a clinical setting.

\section{Conclusion}
A five-block CNN and a feature-fusion hybrid were trained and evaluated for chest X-ray classification. The hybrid improved balanced accuracy relative to the CNN head and CNN-feature-only logistic regression in this single run, while the handcrafted-only model achieved a similar score. High pneumonia recall came with substantial false positives. The results show the importance of evaluating both classes and comparing a hybrid with its individual feature components.

\section{Project Repository}
\url{https://github.com/Ravichandra531/Chest-Xray-pnuemonia}
\end{document}
